
This work investigated the spectral dynamics of SSM adaptation under projection-only LoRA. Three main findings are reported:

\begin{enumerate}
\item The spectral memory horizon H\_spec = -1/log|lambda\_max| is a stable, input-independent property of pretrained Mamba models (CV = 2.22 x $10^-16$ across 1000 sequences), and it predicts perplexity degradation on real text (degradation ratio = 3.03 when context < H\_spec).
\end{enumerate}

\begin{enumerate}
\item Projection-only LoRA achieves exact eigenvalue preservation (delta H\_spec = 0.0\%, eigenvalue correlation = 1.0), confirming that projection parameters and SSM core parameters are architecturally isolated.
\end{enumerate}

\begin{enumerate}
\item Energy redistribution toward slow eigenmodes does not occur (delta E = 5.93 x $10^-7$ nats, six orders of magnitude below threshold). Only 2 of 48 layers in Mamba-1.4B contain slow eigenmodes, with total slow-mode energy at 0.002\%.
\end{enumerate}

The third finding, a negative result, eliminates the Eigenmode Utilization Hypothesis and provides evidence supporting the Memory Horizon Separation Hypothesis. However, direct verification of task failure beyond H\_spec through the planned H-M4 experiment was not completed.

\subsubsection{Future Work}

\begin{itemize}
\item \textbf{Spectral Surgery Methods:} Develop PEFT methods targeting discretization parameters (Delta, A\_log) to test whether eigenvalue modification can extend H\_spec.
\item \textbf{Layer-Selective Adaptation:} Apply SSM-core adaptation selectively to slow-mode layers (18-19) to test whether targeted modification enables energy redistribution.
\item \textbf{Controlled Task Evaluation:} Implement MQAR evaluation with explicit dependency lengths L = {H\_spec, 2\textit{H\_spec, 4}H\_spec} to directly verify task failure at the spectral boundary.
\item \textbf{Cross-Architecture Validation:} Validate the spectral analysis framework on RWKV, RetNet, and Mamba-2.
\end{itemize}
